% Generated from project outputs. Do not edit by hand.
\begin{table}[!htbp]
\centering
\begin{threeparttable}
\caption{Including in-kind by subsample: 2016-2024}
\label{tab:descriptive-2016-2024-subsamples-y2}
\small
\setlength{\tabcolsep}{4pt}
\renewcommand{\arraystretch}{1.15}
\begin{tabularx}{\linewidth}{@{}Xrrrrr@{}}
\toprule
Sample & N & Mean t & SD t & Mean t+1 & SD t+1 \\
\midrule
General & 552 & 2,780.6 & 1,673.9 & 2,973.6 & 1,692.8 \\
Men & 552 & 3,117.5 & 1,940.7 & 3,336.1 & 1,938.1 \\
Women & 552 & 2,227.4 & 1,482.3 & 2,412.6 & 1,569.9 \\
Urban & 552 & 3,027.9 & 1,717.1 & 3,205.0 & 1,717.4 \\
Rural & 547 & 2,278.1 & 1,553.9 & 2,503.2 & 1,694.9 \\
Indigenous & 547 & 2,256.1 & 1,581.8 & 2,528.4 & 1,668.2 \\
Non-Indigenous & 460 & 2,955.7 & 1,758.2 & 3,084.1 & 1,722.5 \\
Bottom 99 percent & 552 & 2,573.7 & 1,327.4 & 2,757.2 & 1,346.8 \\
Bottom 90 percent & 552 & 2,122.3 & 914.5 & 2,267.6 & 913.0 \\
Bottom 50 percent & 548 & 1,162.2 & 463.2 & 1,247.3 & 459.9 \\
\bottomrule
\end{tabularx}
\begin{tablenotes}[flushleft]
\footnotesize
\item Monthly income in quetzales. Unweighted descriptive statistics of survey-weighted cohort means. N counts cohort-transition rows, not individuals or independent cohorts. SD is dispersion across cohort means, not the standard error of a mean. Six annual transitions are pooled; the 2019-2021 gap is excluded. Both incomes must be observed. Subsamples overlap; bottom-income groups are cumulative.
\end{tablenotes}
\end{threeparttable}
\end{table}
